% Source-faithful assembly: Mumford, Abelian Varieties, Chapter IV,
% Horizon transcription pages 0247--0250.

\section{Analytic comparison of theta groups and pairings}
\label{mumford-IV24-analytic-comparison-source}
\dcref{M:ch4:IV.24}

\begin{theorem}[Analytic theta-group comparison]
\label{mumford-thm-analytic-theta-commutator}
\dcref{M:ch4:IV.24}
Let $L=L(H,\alpha)$ be a line bundle on $X=V/U$, let $E={\rm Im}\,H$, and
let $\pi:V\to X$ be the natural map. For $u,v\in U^\perp$,
\[
 e^{-2\pi iE(u,v)}=e^L(\pi u,\pi v).
\]
\source{mumford-abelian-varieties:page-0247}
\source{mumford-abelian-varieties:page-0248}
\uses{mumford-def-line-bundle-commutator,mumford-lem-theta-orthogonal-subgroup}
\end{theorem}
\begin{proof}
The analytic extension has composition
\[
 \psi_{\sigma,w}\psi_{\tau,v}=\psi_{\rho,w+v},
 \qquad \rho=\sigma\tau e^{\pi H(v,w)},
\]
and hence
$1\to\mathbb C^*\to\widetilde{\mathcal G}\to V\to1$.
The datum $\alpha$ lifts $U$ by
$i(u)=\psi_{\sigma,u}$ with $\sigma=\alpha(u)e^{2\pi H(u,u)}$; thus
$L(H,\alpha)=(\mathbb C^*\times V)/i(U)$. Its commutator is
\[
 \psi_{\sigma,w}\psi_{\tau,v}\psi_{\sigma,w}^{-1}\psi_{\tau,v}^{-1}
 =\psi_{\widetilde e(v,w),0},
 \qquad \widetilde e(v,w)=e^{2\pi iE(v,w)}.
\]
Therefore $p^{-1}(U^\perp)$ is the centralizer of $i(U)$, and its elements
descend to automorphisms of $L(H,\alpha)$. Since
$K(L(H,\alpha))\simeq U^\perp/U$, the extensions
\[
1\to\mathbb C^*\to\widetilde{\mathcal G}_0/i(U)\to U^\perp/U\to1,
\qquad
1\to\mathbb C^*\to\mathcal G(L(H,\alpha))\to K(L(H,\alpha))\to1
\]
are isomorphic. Their commutators agree, giving the formula.
\end{proof}

\begin{proposition}[Analytic and l-adic Riemann pairings]
\label{mumford-prop-analytic-ladic-pairing}
\dcref{M:ch4:IV.24}
For $X=V/U$ and $\ell$, the canonical pairing $E^L$ on $T_\ell X$ is, up to
sign, the $\mathbb Z_\ell$-linear extension of $E:U\times U\to\mathbb Z$.
In particular, for the Poincare bundle on $Y\times\widehat Y$, the analytic
pairing of $U,\widehat U$ agrees up to sign with the canonical pairing of
$T_\ell Y,T_\ell\widehat Y$.
\source{mumford-abelian-varieties:page-0248}
\source{mumford-abelian-varieties:page-0249}
\uses{mumford-thm-analytic-theta-commutator}
\end{proposition}
\begin{proof}
Let $M_\ell=\varprojlim_n\mu_{\ell^n}$ with basis
$\zeta=(e^{2\pi i/\ell^n})_n$. If $\pi_\ell(u)$ is represented by
$u_n=\pi(u/\ell^n)$, then
\[
\begin{aligned}
 E^L(\pi_\ell u,\pi_\ell v)
 &=\bigl(e_{\ell^n}^L(u_n,v_n)\bigr)_n
 =\bigl(e^{L^{\ell^n}}(u_n,v_n)\bigr)_n\\
 &=\bigl(e^{-2\pi i\ell^nE(u/\ell^n,v/\ell^n)}\bigr)_n
 =\bigl((\zeta_{\ell^n})^{-E(u,v)}\bigr)_n
 =-E(u,v)\cdot\zeta.
\end{aligned}
\]
The Poincare-bundle assertion is the corresponding non-degenerate case.
\end{proof}

\begin{proposition}[Analytic interpretation and positivity of Rosati]
\label{mumford-prop-analytic-rosati-positivity}
\dcref{M:ch4:IV.24}
Analytically,
\[
{\rm End}^0(X)=\{T:V\to V\text{ complex-linear}:T(\mathbb Q\cdot U)\subset\mathbb Q\cdot U\},
\]
and the Rosati involution is the $H$-adjoint $T^*$ defined by
$H(T^*x,y)=H(x,Ty)$. It equals the algebraic Rosati involution and satisfies
\[
{\rm Tr}(T^*T)\ge0,\qquad T\ne0.
\]
\source{mumford-abelian-varieties:page-0249}
\source{mumford-abelian-varieties:page-0250}
\uses{mumford-prop-analytic-ladic-pairing}
\end{proposition}
\begin{proof}
If $x\in\mathbb Q\cdot U$ and $y\in\mathbb Q\cdot U$, then
$E(T^*x,y)=E(x,Ty)\in\mathbb Q$, so $T^*$ preserves $\mathbb Q\cdot U$.
For $T'$ algebraic Rosati and $x,y\in U$,
\[
\begin{aligned}
 E((T^*-T')x,y)&={\rm Im}\,H(T^*x,y)+E^L(T'\pi x,T\pi y)\\
 &={\rm Im}\,H(x,Ty)+E^L(\pi x,T\pi y)=0,
\end{aligned}
\]
so $T^*=T'$. The operator $T^*T$ is positive self-adjoint for the positive
Hermitian form $H$, so its complex trace is positive. Its trace on the rational
space is twice the complex trace, and its $\ell$-adic trace on
$T_\ell X\simeq U\otimes\mathbb Z_\ell$ is the same rational trace. Hence
positivity follows.
\end{proof}
